\documentclass[11pt]{letter}
\usepackage[margin=1in]{geometry}
\usepackage{parskip}

\signature{Hikaru Wakaura \\ Taiki Tanimae \\ QIRI (Quantum Integrated
Research Institute Inc.)}
\address{QIRI (Quantum Integrated Research Institute Inc.)\\
1--16--3 Akasaka, Minato-ku\\ Tokyo 107--0061, Japan\\
h.wakaura@deeptell.jp,\ t.tanimae@deeptell.jp}

\begin{document}
\begin{letter}{Editorial Office \\ Physical Review A \\ American Physical
Society}

\opening{Dear Editors,}

We submit the accompanying manuscript, ``Which quantum algorithms suit
photonic continuous-variable hardware? A comparative resource and
noise-robustness study,'' for consideration as a Regular Article in
\emph{Physical Review A}. The work falls squarely within the journal's
scope of quantum information, quantum computation, and the physics of
quantum systems.

\textbf{Motivation.} Photonic continuous-variable (CV) architectures
support deterministic, room-temperature generation of large cluster
states and are a leading route to fault-tolerant quantum computing,
where overhead is dominated by squeezing and by the consumption of
non-Gaussian resource states. Different algorithms place very different
demands on these resources, yet no unified comparison has quantified
which algorithm classes are intrinsically robust on CV hardware and how
much Gottesman--Kitaev--Preskill (GKP) error correction actually helps.

\textbf{Main results.} We port four representative algorithms---%
control-free phase estimation, randomized product formulas, quantum
Kolmogorov--Arnold networks, and Regev factoring---to a common CV model
and compare them under photon loss (exact Fock-space Kraus channel) and
Gaussian displacement with GKP correction. We establish: (i)~control-%
free phase estimation is the most noise-robust because it needs no
controlled unitaries and, for quadratic Hamiltonians, no non-Gaussian
gates, a property that survives for non-Gaussian Hamiltonians;
(ii)~GKP correction reduces residual error only above the finite-%
squeezing floor (proved); whether an algorithm benefits is set by a
noise- versus systematic-limited dichotomy whose empirical envelope
constants we report with $R^{2}\!\ge\!0.89$; (iii)~the comb readout of
Regev requires squeezing scaling as $\sqrt{n}$ (proved), demonstrating
that CV-native cheapness of state preparation and the Fourier transform
does not rescue a depth- and readout-heavy algorithm; (iv)~against a
qubit surface-code baseline, the CV advantage is structural---Gaussian
operations are free---rather than threshold-based.

\textbf{Methodological rigor.} The core GKP-correction surrogate is
validated against an explicit Glancy--Knill error-correction Monte Carlo
to within $0.1\%$, lifting a common concern in CV-resource analyses. A
two-dimensional noise phase diagram and Hamiltonian ensembles confirm
the predictions are not artefacts of single operating points or single
instances. We are explicit about which results are exact, which are
validated model surrogates, and which are extrapolations.

\textbf{New worked physics examples.} Responding to earlier feedback that
the work lacked concrete physics results, the manuscript now includes
quantitative demonstrations on \emph{three diatomic molecules}---H$_2$,
HF, and I$_2$, spanning a decade in vibrational quantum
($\omega_e=214$--$4401$~cm$^{-1}$) and an order of magnitude in
anharmonicity ($x_e=0.003$--$0.028$)---using experimentally fixed
Huber--Herzberg Morse parameters (Sec.~VIII). CV control-free QPE
recovers the populated vibrational eigenenergies of all three molecules
to within $0.55$~cm$^{-1}$ \emph{below the standard $1$~cm$^{-1}$
spectroscopic threshold}; weakly anharmonic I$_2$ reaches
$\sim\!10^{-2}$~cm$^{-1}$. We further show that matching this accuracy
with a discrete-variable QPE requires $\sim 1.5\times10^{10}$ T~gates and
$\sim\!16{,}000$ controlled time evolutions, whereas the CV control-free
protocol needs zero controlled unitaries and zero non-Gaussian gates for
the harmonic part. This is, to our knowledge, the first end-to-end
CV-resource-quantified calculation of \emph{multiple} real molecular
spectra at spectroscopic accuracy, and we believe it constitutes a
meaningful physics contribution beyond methodology. A Chebyshev-KAN
classifier trained on the recovered eigenstates assigns the vibrational
quantum number with $100\%$ accuracy under realistic measurement noise.

\textbf{Submission history.} An earlier short-form version of this work
was returned by \emph{Applied Physics Express} on scope grounds (the
referee judged the topic to lie in fundamental quantum-information
physics rather than applied photonics) and a preliminary version of the
present manuscript was returned by PRA on the grounds that the original
content did not contain significant new physics results. The H$_2$
vibrational-spectrum calculation reported in Sec.~VIII directly
addresses that concern, providing a concrete, quantitative physics
deliverable while retaining the original cross-algorithm methodology. The manuscript has not been
published elsewhere and is not under consideration at any other journal;
the authors declare no competing interests. All scripts and numerical
results that reproduce every figure and table are available with the
submission.

We thank you for your consideration.

\closing{Sincerely,}

\end{letter}
\end{document}
